% Generisano: uređuj beleske/sNNN.tex i naslove u manifestu.
\documentclass[11pt,a4paper]{article}
\input{../_alati/beleske-v1.tex}
\newcommand{\BeleskePrezentacija}{build/03_kola_srednjeg_stepena_integracije.pdf}
\begin{document}
\BeleskaSlajda{03-s001}{1}{Kola srednjeg stepena integracije}{beleske/s001.tex}
\BeleskaSlajda{03-s002}{2}{Stepeni integracije}{beleske/s002.tex}
\BeleskaSlajda{03-s003}{3}{Pakovanja i priključci integrisanih kola}{beleske/s003.tex}
\BeleskaSlajda{03-s004}{4}{Izbor kola i neiskorišćeni ulazi}{beleske/s004.tex}
\BeleskaSlajda{03-s005}{5}{Potpuni binarni dekoder 2/4}{beleske/s005.tex}
\BeleskaSlajda{03-s006}{6}{Primer funkcije sa dekoderom}{beleske/s006.tex}
\BeleskaSlajda{03-s007}{7}{Simbol dekodera i oznake signala}{beleske/s007.tex}
\BeleskaSlajda{03-s008}{8}{Izbor komponente signalom CS}{beleske/s008.tex}
\BeleskaSlajda{03-s009}{9}{Piramidalna struktura dekodera}{beleske/s009.tex}
\BeleskaSlajda{03-s010}{10}{Matrična struktura: povezivanje dekodera}{beleske/s010.tex}
\BeleskaSlajda{03-s010b}{11}{Matrična i piramidalna struktura: poređenje}{beleske/s010b.tex}
\BeleskaSlajda{03-s011}{12}{Aktivni nivoi i negacije}{beleske/s011.tex}
\BeleskaSlajda{03-s012}{13}{Demultipleksor 1/4}{beleske/s012.tex}
\BeleskaSlajda{03-s013}{14}{Multiplekser 4/1}{beleske/s013.tex}
\BeleskaSlajda{03-s014}{15}{Realizacija funkcije multiplekserom}{beleske/s014.tex}
\BeleskaSlajda{03-s015}{16}{Izbor i proširenje multipleksera}{beleske/s015.tex}
\BeleskaSlajda{03-s016}{17}{Koder prioriteta: potreba za signalom GS}{beleske/s016.tex}
\BeleskaSlajda{03-s017}{18}{Koder prioriteta sa izlazom GS}{beleske/s017.tex}
\BeleskaSlajda{03-s018}{19}{Kompaktnija realizacija kodera prioriteta}{beleske/s018.tex}
\BeleskaSlajda{03-s019}{20}{Koder prioriteta sa 16 ulaza}{beleske/s019.tex}
\BeleskaSlajda{03-s020}{21}{Signali EI i EO u koderu prioriteta}{beleske/s020.tex}
\BeleskaSlajda{03-s021}{22}{Lanac kodera prioriteta sa 16 ulaza}{beleske/s021.tex}
\end{document}
